\chapter{1955:Vincent du Vigneaud}

\label{chap:chemistry1955}

The Nobel Prize in Chemistry for the year 1955 was awarded to Vincent du Vigneaud.

\textbf{Motivation:}

for his work on biochemically important sulphur compounds, especially for the first synthesis of a polypeptide hormone

\section{Vincent du Vigneaud}

\subsection*{Early Life and Education}

Vincent du Vigneaud was born on 1901-05-18 in Chicago, Illinois, USA.

\subsection*{Career and Major Research}

du Vigneaud studied at the University of Illinois and received his doctorate from the University of Rochester in 1927. He worked at the University of Chicago and Cornell University Medical College in New York, where he became professor of biochemistry. He studied the sulfur-containing amino acids methionine, cysteine, and cystine, and investigated the sulfur chemistry of the hormones insulin, vasopressin, and oxytocin. In 1953 his group achieved the first synthesis of the peptide hormone oxytocin, and in 1954 the synthesis of vasopressin.

\subsection*{Nobel Prize-winning Work}

The work for which Vincent du Vigneaud was awarded the Nobel Prize in Chemistry in 1955 was described by the Nobel Committee as: "for his work on biochemically important sulphur compounds, especially for the first synthesis of a polypeptide hormone".

du Vigneaud determined the sequence of the eight amino acids in oxytocin and vasopressin, including the disulfide bridge formed by their cystine residues, and then built the complete hormone chains. The synthetic oxytocin was shown to possess the full biological activity of the natural hormone, proving the structure correct.

\subsection*{Significance and Impact}

The synthesis of oxytocin was the first total synthesis of a peptide hormone and demonstrated that a synthetic polypeptide could reproduce the activity of a natural hormone. It established the chemistry of the sulfur-containing amino acids and opened the way to the synthesis and medical use of peptide hormones.

\subsection*{Later Career and Honors}

du Vigneaud remained professor of biochemistry at Cornell University Medical College until his retirement, continuing research on peptide and hormone chemistry. He died on 1978-12-11 in White Plains, New York.

\subsection*{Legacy}

The oxytocin synthesis is regarded as the founding achievement of peptide-hormone chemistry, and du Vigneaud's methods anticipated the field of peptide synthesis that later produced insulin, growth hormone, and a wide range of therapeutic peptides.

\subsection*{Selected Publications}

\begin{itemize}

\item \emph{The Synthesis of an Octapeptide Amide with the Hormonal Activity of Oxytocin}, Journal of the American Chemical Society, 1953.

\end{itemize}

\clearpage

\section*{Chapter Summary}

The 1955 prize recognized the first synthesis of a polypeptide hormone. Vincent du Vigneaud determined the structure of oxytocin and vasopressin and synthesized them in his laboratory at Cornell, work built on his long study of the biochemically important sulfur-containing amino acids.

